\title{2022 年高考数学 全国甲卷-理科}
\date{}
\maketitle

\section*{一、选择题}
本题共12小题，每小题5分，共60分.在每小题给出的四个选项中，只有一项是符合题目要求的.

\begin{question}[points=5]
若 $z = -1 + \sqrt{3}i$，则 $\frac{z}{z\bar{z} - 1} = $（ ）
\begin{choices}
  \item $-1+\sqrt{3}i$
  \item $-1-\sqrt{3}i$
  \item $-\frac{1}{3}+\frac{\sqrt{3}}{3}i$
  \item $-\frac{1}{3}-\frac{\sqrt{3}}{3}i$
\end{choices}
\begin{solution}
\textbf{答案：}C

\textbf{分析：}由共轭复数的概念及复数的运算即可得解.

\textbf{详解：}$\bar{z} = -1 - \sqrt{3}i$, $z\bar{z} = (-1+\sqrt{3}i)(-1-\sqrt{3}i) = 1+3 = 4$. 则 $\frac{z}{z\bar{z}-1} = \frac{-1+\sqrt{3}i}{3} = -\frac{1}{3} + \frac{\sqrt{3}}{3}i$. 故选C.
\end{solution}
\end{question}

\begin{question}[points=5]
某社区通过公益讲座以普及社区居民的垃圾分类知识.为了解讲座效果，随机抽取10位社区居民，让他们在讲座前和讲座后各回答一份垃圾分类知识问卷，这10位社区居民在讲座前和讲座后问卷答题的正确率如下图：



则（ ）

\begin{center}
  \includegraphics[width=.38\linewidth]{question-2-1.jpg}
\end{center}
\begin{choices}
  \item 讲座前问卷答题的正确率的中位数小于$70\%$
  \item 讲座后问卷答题的正确率的平均数大于$85\%$
  \item 讲座前问卷答题的正确率的标准差小于讲座后正确率的标准差
  \item 讲座后问卷答题的正确率的极差大于讲座前正确率的极差
\end{choices}
\begin{solution}
\textbf{答案：}B

\textbf{分析：}由图表信息，结合中位数、平均数、标准差、极差的概念，逐项判断即可得解.

\textbf{详解：}讲座前中位数为 $\frac{70\%+75\%}{2} > 70\%$，所以A错；讲座后问卷答题的正确率只有一个是$80\%$，4个$85\%$，剩下全部大于等于$90\%$，所以讲座后问卷答题的正确率的平均数大于$85\%$，所以B对；讲座前问卷答题的正确率更加分散，所以讲座前问卷答题的正确率的标准差大于讲座后正确率的标准差，所以C错；讲座后问卷答题的正确率的极差为$100\%-80\%=20\%$，讲座前问卷答题的正确率的极差为$95\%-60\%=35\%>20\%$，所以D错.故选B.
\end{solution}
\end{question}

\begin{question}[points=5]
设全集 $U=\{-2,-1,0,1,2,3\}$，集合 $A=\{-1,2\}$，$B=\{x \mid x^2-4x+3=0\}$，则 $\complement_U(A \cup B)=$（ ）
\begin{choices}
  \item $\{1,3\}$
  \item $\{0,3\}$
  \item $\{-2,1\}$
  \item $\{-2,0\}$
\end{choices}
\begin{solution}
\textbf{答案：}D

\textbf{分析：}解方程求出集合B，再由集合的运算即可得解.

\textbf{详解：}由题意，$B=\{x\mid x^2-4x+3=0\}=\{1,3\}$，所以$A\cup B=\{-1,1,2,3\}$，所以$\complement_U(A\cup B)=\{-2,0\}$. 故选D.
\end{solution}
\end{question}

\begin{question}[points=5]
如图，网格纸上绘制的是一个多面体的三视图，网格小正方形的边长为1，则该多面体的体积为（ ）

\begin{center}
  \includegraphics[width=.38\linewidth]{question-4-1.jpg}
\end{center}
\begin{choices}
  \item $8$
  \item $12$
  \item $16$
  \item $20$
\end{choices}
\begin{solution}
\textbf{答案：}B

\textbf{分析：}由三视图还原几何体，再由棱柱的体积公式即可得解.

\textbf{详解：}由三视图还原几何体如图，则该直四棱柱的体积 $V = \frac{2+4}{2} \times 2 \times 2 = 12$. 故选B.
\end{solution}
\end{question}

\begin{question}[points=5]
函数 $y=(3^x-3^{-x})\cos x$ 在区间 $\left[-\frac{\pi}{2}, \frac{\pi}{2}\right]$ 的图象大致为（ ）

\begin{center}
  \includegraphics[width=.38\linewidth]{question-5-1.jpg}
\end{center}
\begin{choices}
  \item A选项图
  \item B选项图
  \item C选项图
  \item D选项图
\end{choices}
\begin{solution}
\textbf{答案：}A

\textbf{分析：}由函数的奇偶性结合指数函数、三角函数的性质逐项排除即可得解.

\textbf{详解：}令 $f(x)=(3^x-3^{-x})\cos x$, $x\in\left[-\frac{\pi}{2},\frac{\pi}{2}\right]$, 则 $f(-x)=(3^{-x}-3^x)\cos(-x)=-(3^x-3^{-x})\cos x=-f(x)$, 所以 $f(x)$ 为奇函数，排除BD；又当 $x\in\left(0,\frac{\pi}{2}\right)$ 时，$3^x-3^{-x}>0$, $\cos x>0$, 所以 $f(x)>0$, 排除C. 故选A.
\end{solution}
\end{question}

\begin{question}[points=5]
当 $x=1$ 时，函数 $f(x)=a\ln x + \frac{b}{x}$ 取得最大值 $-2$，则 $f'(2)=$（ ）
\begin{choices}
  \item $-1$
  \item $-\frac{1}{2}$
  \item $\frac{1}{2}$
  \item $1$
\end{choices}
\begin{solution}
\textbf{答案：}B

\textbf{分析：}根据题意可知 $f(1)=-2$, $f'(1)=0$ 即可解得 $a,b$，再根据 $f'(x)$ 即可解出.

\textbf{详解：}因为函数 $f(x)$ 定义域为 $(0,+\infty)$，所以依题可知 $f(1)=-2$, $f'(1)=0$，而 $f'(x)=\frac{a}{x}-\frac{b}{x^2}$，所以 $b=-2$, $a-b=0$，即 $a=-2$, $b=-2$，所以 $f'(x)=-\frac{2}{x}+\frac{2}{x^2}$，因此函数 $f(x)$ 在 $(0,1)$ 上递增，在 $(1,+\infty)$ 上递减，$x=1$ 时取最大值，满足题意，即有 $f'(2)=-1+\frac{1}{2}=-\frac{1}{2}$. 故选B.
\end{solution}
\end{question}

\begin{question}[points=5]
在长方体 $ABCD-A_1B_1C_1D_1$ 中，已知 $B_1D$ 与平面 $ABCD$ 和平面 $AA_1B_1B$ 所成的角均为 $30^\circ$，则（ ）
\begin{choices}
  \item $AB=2AD$
  \item $AB$ 与平面 $AB_1C_1D$ 所成的角为 $30^\circ$
  \item $AC=CB_1$
  \item $B_1D$ 与平面 $BB_1C_1C$ 所成的角为 $45^\circ$
\end{choices}
\begin{solution}
\textbf{答案：}D

\textbf{分析：}根据线面角的定义以及长方体的结构特征即可求出.

\textbf{详解：}不妨设 $AB=a$, $AD=b$, $AA_1=c$，依题以及长方体的结构特征可知，$B_1D$ 与平面 $ABCD$ 所成角为 $\angle B_1DB$，$B_1D$ 与平面 $AA_1B_1B$ 所成角为 $\angle DB_1A$，所以 $\sin30^\circ=\frac{c}{B_1D}=\frac{b}{B_1D}$，即 $b=c$, $B_1D=2c=\sqrt{a^2+b^2+c^2}$，解得 $a=\sqrt{2}c$. 对于A，$AB=a$, $AD=b$, $AB=\sqrt{2}AD$，A错误；对于B，过B作 $BE\perp AB_1$ 于E，易知 $BE\perp$ 平面 $AB_1C_1D$，所以AB与平面 $AB_1C_1D$ 所成角为 $\angle BAE$，因为 $\tan\angle BAE=\frac{c}{a}=\frac{\sqrt{2}}{2}$，所以 $\angle BAE\neq30^\circ$，B错误；对于C，$AC=\sqrt{a^2+b^2}=\sqrt{3}c$, $CB_1=\sqrt{b^2+c^2}=\sqrt{2}c$，$AC\neq CB_1$，C错误；对于D，$B_1D$ 与平面 $BB_1C_1C$ 所成角为 $\angle DB_1C$，$\sin\angle DB_1C=\frac{CD}{B_1D}=\frac{a}{2c}=\frac{\sqrt{2}}{2}$，而 $0<\angle DB_1C<90^\circ$，所以 $\angle DB_1C=45^\circ$，D正确. 故选D.
\end{solution}
\end{question}

\begin{question}[points=5]
沈括的《梦溪笔谈》是中国古代科技史上的杰作，其中收录了计算圆弧长度的“会圆术”，如图，$\overset{\frown}{AB}$ 是以O为圆心，OA为半径的圆弧，C是AB的中点，D在 $\overset{\frown}{AB}$ 上，$CD\perp AB$. “会圆术”给出 $\overset{\frown}{AB}$ 的弧长近似值s的计算公式：$s=AB+\frac{CD^2}{OA}$. 当 $OA=2$, $\angle AOB=60^\circ$ 时，$s=$（ ）

\begin{center}
  \includegraphics[width=.38\linewidth]{question-8-1.jpg}
\end{center}
\begin{choices}
  \item $\frac{11-3\sqrt{3}}{2}$
  \item $\frac{11-4\sqrt{3}}{2}$
  \item $\frac{9-3\sqrt{3}}{2}$
  \item $\frac{9-4\sqrt{3}}{2}$
\end{choices}
\begin{solution}
\textbf{答案：}B

\textbf{分析：}连接OC，分别求出AB, OC, CD，再根据题中公式即可得出答案.

\textbf{详解：}连接OC，因为C是AB的中点，所以 $OC\perp AB$，又 $CD\perp AB$，所以 $O,C,D$ 三点共线，即 $OD=OA=OB=2$，又 $\angle AOB=60^\circ$，所以 $AB=OA=OB=2$，则 $OC=\sqrt{3}$，故 $CD=2-\sqrt{3}$，所以 $s=AB+\frac{CD^2}{OA}=2+\frac{(2-\sqrt{3})^2}{2}=2+\frac{4-4\sqrt{3}+3}{2}=2+\frac{7-4\sqrt{3}}{2}=\frac{11-4\sqrt{3}}{2}$. 故选B.
\end{solution}
\end{question}

\begin{question}[points=5]
甲、乙两个圆锥的母线长相等，侧面展开图的圆心角之和为 $2\pi$，侧面积分别为 $S_\text{甲}$ 和 $S_\text{乙}$，体积分别为 $V_\text{甲}$ 和 $V_\text{乙}$. 若 $\frac{S_\text{甲}}{S_\text{乙}}=2$，则 $\frac{V_\text{甲}}{V_\text{乙}}=$（ ）
\begin{choices}
  \item $\sqrt{5}$
  \item $2\sqrt{2}$
  \item $\sqrt{10}$
  \item $\frac{5\sqrt{10}}{4}$
\end{choices}
\begin{solution}
\textbf{答案：}C

\textbf{分析：}设母线长为l，甲圆锥底面半径为$r_1$，乙圆锥底面圆半径为$r_2$，根据圆锥的侧面积公式可得$r_1=2r_2$，再结合圆心角之和可将$r_1,r_2$分别用l表示，再利用勾股定理分别求出两圆锥的高，再根据圆锥的体积公式即可得解.

\textbf{详解：}设母线长为l，甲圆锥底面半径为$r_1$，乙圆锥底面圆半径为$r_2$，则 $\frac{S_\text{甲}}{S_\text{乙}}=\frac{\pi r_1 l}{\pi r_2 l}=\frac{r_1}{r_2}=2$，所以 $r_1=2r_2$，又 $\frac{2\pi r_1}{l}+\frac{2\pi r_2}{l}=2\pi$，则 $\frac{r_1+r_2}{l}=1$，所以 $r_1=\frac{2}{3}l$, $r_2=\frac{1}{3}l$，所以甲圆锥的高 $h_1=\sqrt{l^2-\frac{4}{9}l^2}=\frac{\sqrt{5}}{3}l$，乙圆锥的高 $h_2=\sqrt{l^2-\frac{1}{9}l^2}=\frac{2\sqrt{2}}{3}l$，所以 $\frac{V_\text{甲}}{V_\text{乙}}=\frac{\frac{1}{3}\pi r_1^2 h_1}{\frac{1}{3}\pi r_2^2 h_2}=\frac{\frac{4}{9}l^2 \times \frac{\sqrt{5}}{3}l}{\frac{1}{9}l^2 \times \frac{2\sqrt{2}}{3}l}=\sqrt{10}$. 故选C.
\end{solution}
\end{question}

\begin{question}[points=5]
椭圆 $C:\frac{x^2}{a^2}+\frac{y^2}{b^2}=1(a>b>0)$ 的左顶点为A，点P，Q均在C上，且关于y轴对称. 若直线AP, AQ的斜率之积为 $\frac{1}{4}$，则C的离心率为（ ）
\begin{choices}
  \item $\frac{\sqrt{3}}{2}$
  \item $\frac{\sqrt{2}}{2}$
  \item $\frac{1}{2}$
  \item $\frac{1}{3}$
\end{choices}
\begin{solution}
\textbf{答案：}A

\textbf{分析：}设$P(x_1,y_1)$，则$Q(-x_1,y_1)$，根据斜率公式结合题意可得$\frac{y_1^2}{-x_1^2+a^2}=\frac{1}{4}$，再根据$\frac{x_1^2}{a^2}+\frac{y_1^2}{b^2}=1$，将$y_1$用$x_1$表示，整理，再结合离心率公式即可得解.

\textbf{详解：}设$P(x_1,y_1)$，则$Q(-x_1,y_1)$，则由$k_{AP}\cdot k_{AQ}=\frac{1}{4}$得：$k_{AP}\cdot k_{AQ}=\frac{y_1}{x_1+a}\cdot\frac{y_1}{-x_1+a}=\frac{y_1^2}{-x_1^2+a^2}=\frac{1}{4}$. 由$\frac{x_1^2}{a^2}+\frac{y_1^2}{b^2}=1$得$y_1^2=\frac{b^2(a^2-x_1^2)}{a^2}$，代入得$\frac{b^2(a^2-x_1^2)}{a^2(-x_1^2+a^2)}=\frac{1}{4}$，即$\frac{b^2}{a^2}=\frac{1}{4}$，所以椭圆C的离心率$e=\frac{c}{a}=\sqrt{1-\frac{b^2}{a^2}}=\sqrt{1-\frac{1}{4}}=\frac{\sqrt{3}}{2}$. 故选A.
\end{solution}
\end{question}

\begin{question}[points=5]
设函数 $f(x)=\sin\left(\omega x+\frac{\pi}{3}\right)$ 在区间 $(0,\pi)$ 恰有三个极值点、两个零点，则 $\omega$ 的取值范围是（ ）
\begin{choices}
  \item $\left[\frac{5}{3},\frac{13}{6}\right)$
  \item $\left[\frac{5}{3},\frac{19}{6}\right)$
  \item $\left(\frac{13}{6},\frac{8}{3}\right]$
  \item $\left(\frac{13}{6},\frac{19}{6}\right]$
\end{choices}
\begin{solution}
\textbf{答案：}C

\textbf{分析：}由x的取值范围得到$\omega x+\frac{\pi}{3}$的取值范围，再结合正弦函数的性质得到不等式组，解得即可.

\textbf{详解：}依题意可得$\omega>0$，因为$x\in(0,\pi)$，所以$\omega x+\frac{\pi}{3}\in\left(\frac{\pi}{3},\omega\pi+\frac{\pi}{3}\right)$，要使函数在区间$(0,\pi)$恰有三个极值点、两个零点，又$y=\sin x$, $x\in\left(\frac{\pi}{3},3\pi\right)$的图象如图所示，则$\frac{5\pi}{2}<\omega\pi+\frac{\pi}{3}\leq 3\pi$，解得$\frac{13}{6}<\omega\leq\frac{8}{3}$，即$\omega\in\left(\frac{13}{6},\frac{8}{3}\right]$. 故选C.
\end{solution}
\end{question}

\begin{question}[points=5]
已知 $a=\frac{31}{32}$, $b=\cos\frac{1}{4}$, $c=4\sin\frac{1}{4}$，则（ ）
\begin{choices}
  \item $c>b>a$
  \item $b>a>c$
  \item $a>b>c$
  \item $a>c>b$
\end{choices}
\begin{solution}
\textbf{答案：}A

\textbf{分析：}由$\frac{c}{b}=4\tan\frac{1}{4}$结合三角函数的性质可得$c>b$；构造函数$f(x)=\cos x+\frac{1}{2}x^2-1, x\in(0,+\infty)$，利用导数可得$b>a$，即可得解.

\textbf{详解：}因为当$x\in(0,\frac{\pi}{2})$时，$x<\tan x$，故$\frac{c}{b}=4\tan\frac{1}{4}>1$，所以$c>b$；设$f(x)=\cos x+\frac{1}{2}x^2-1$, $x\in(0,+\infty)$，$f'(x)=-\sin x+x>0$，所以$f(x)$在$(0,+\infty)$单调递增，故$f(\frac{1}{4})>f(0)=0$，所以$\cos\frac{1}{4}-\frac{31}{32}>0$，所以$b>a$，因此$c>b>a$. 故选A.
\end{solution}
\end{question}

\section*{二、填空题}
本题共4小题，每小题5分，共20分.

\begin{question}[points=5]
设向量 $\vec{a}$, $\vec{b}$ 的夹角的余弦值为 $\frac{1}{3}$，且 $|\vec{a}|=1$，$|\vec{b}|=3$，则 $(2\vec{a}+\vec{b})\cdot\vec{b}=$ ．
\fillin[{$11$}]
\begin{solution}
\textbf{答案：}$11$

\textbf{分析：}设$\vec{a}$与$\vec{b}$的夹角为$\theta$，依题意可得$\cos\theta=\frac{1}{3}$，再根据数量积的定义求出$\vec{a}\cdot\vec{b}$，最后根据数量积的运算律计算可得.

\textbf{详解：}解：设$\vec{a}$与$\vec{b}$的夹角为$\theta$，因为$\vec{a}$与$\vec{b}$的夹角的余弦值为$\frac{1}{3}$，即$\cos\theta=\frac{1}{3}$，又$|\vec{a}|=1$, $|\vec{b}|=3$，所以$\vec{a}\cdot\vec{b}=|\vec{a}||\vec{b}|\cos\theta=1\times3\times\frac{1}{3}=1$，所以$(2\vec{a}+\vec{b})\cdot\vec{b}=2\vec{a}\cdot\vec{b}+|\vec{b}|^2=2\times1+3^2=11$. 故答案为11.
\end{solution}
\end{question}

\begin{question}[points=5]
若双曲线 $y^2-\frac{x^2}{m}=1$ ($m>0$) 的渐近线与圆 $x^2+y^2-4y+3=0$ 相切，则 $m=$ ．
\fillin[{$\frac{\sqrt{3}}{3}$}]
\begin{solution}
\textbf{答案：}$\frac{\sqrt{3}}{3}$

\textbf{分析：}首先求出双曲线的渐近线方程，再将圆的方程化为标准式，即可得到圆心坐标与半径，依题意圆心到直线的距离等于圆的半径，即可得到方程，解得即可.

\textbf{详解：}解：双曲线 $y^2-\frac{x^2}{m}=1$ ($m>0$) 的渐近线为 $y=\pm\frac{x}{\sqrt{m}}$，即 $x\pm\sqrt{m}y=0$，不妨取 $x+\sqrt{m}y=0$，圆 $x^2+y^2-4y+3=0$ 即 $x^2+(y-2)^2=1$，所以圆心为 $(0,2)$，半径 $r=1$，依题意圆心 $(0,2)$ 到渐近线 $x+\sqrt{m}y=0$ 的距离 $d=\frac{|2\sqrt{m}|}{\sqrt{1+m}}=1$，解得 $m=\frac{\sqrt{3}}{3}$ 或 $m=-\frac{\sqrt{3}}{3}$（舍去）. 故答案为 $\frac{\sqrt{3}}{3}$.
\end{solution}
\end{question}

\begin{question}[points=5]
从正方体的8个顶点中任选4个，则这4个点在同一个平面的概率为．
\fillin[{$\frac{6}{35}$}]
\begin{solution}
\textbf{答案：}$\frac{6}{35}$

\textbf{分析：}根据古典概型的概率公式即可求出.

\textbf{详解：}从正方体的8个顶点中任取4个，有 $n=C_8^4=70$ 个结果，这4个点在同一个平面的有 $m=6+6=12$ 个，故所求概率 $P=\frac{m}{n}=\frac{12}{70}=\frac{6}{35}$. 故答案为 $\frac{6}{35}$.
\end{solution}
\end{question}

\begin{question}[points=5]
已知 $\triangle ABC$ 中，点D在边BC上，$\angle ADB=120^\circ$, $AD=2$, $CD=2BD$. 当 $\frac{AC}{AB}$ 取得最小值时，$BD=$ ．
\fillin[{$\sqrt{3}-1$}]
\begin{solution}
\textbf{答案：}$\sqrt{3}-1$

\textbf{分析：}设$CD=2BD=2m>0$，利用余弦定理表示出$\frac{AC^2}{AB^2}$后，结合基本不等式即可得解.

\textbf{详解：}设$CD=2BD=2m>0$，则在$\triangle ABD$中，$AB^2=BD^2+AD^2-2BD\cdot AD\cos\angle ADB=m^2+4+2m$，在$\triangle ACD$中，$AC^2=CD^2+AD^2-2CD\cdot AD\cos\angle ADC=4m^2+4-4m$，所以 $\frac{AC^2}{AB^2}=\frac{4m^2+4-4m}{m^2+4+2m}=\frac{4(m^2+4+2m)-12(1+m)}{m^2+4+2m}=4-\frac{12}{m+1+\frac{3}{m+1}} \geq 4-\frac{12}{2\sqrt{3}}=4-2\sqrt{3}$，当且仅当 $m+1=\frac{3}{m+1}$ 即 $m=\sqrt{3}-1$ 时等号成立，所以当 $\frac{AC}{AB}$ 取最小值时，$m=\sqrt{3}-1$，即 $BD=\sqrt{3}-1$. 故答案为 $\sqrt{3}-1$.
\end{solution}
\end{question}

\section*{三、解答题}
共70分.解答应写出文字说明、证明过程或演算步骤.第17～21题为必考题，每个试题考生都必须作答.第22、23题为选考题，考生根据要求作答.

\begin{question}[points=12]
记 $S_n$ 为数列 $\{a_n\}$ 的前n项和.已知 $\frac{2S_n}{n}+n=2a_n+1$.

(1) 证明：$\{a_n\}$ 是等差数列；

(2) 若 $a_4, a_7, a_9$ 成等比数列，求 $S_n$ 的最小值.
\begin{solution}
\textbf{答案：}(1) 证明见解析；(2) $-78$

\textbf{分析：}(1) 依题意可得 $2S_n+n^2=2na_n+n$，根据 $a_n=\begin{cases} S_1, & n=1 \\ S_n-S_{n-1}, & n\ge 2 \end{cases}$，作差即可得到 $a_n-a_{n-1}=1$，从而得证. (2) 由(1)及等比中项的性质求出 $a_1$，即可得到 $\{a_n\}$ 的通项公式与前n项和，再根据二次函数的性质计算可得.

\textbf{详解：}(1) 因为 $\frac{2S_n}{n}+n=2a_n+1$，即 $2S_n+n^2=2na_n+n$ ①，当 $n\ge 2$ 时，$2S_{n-1}+(n-1)^2=2(n-1)a_{n-1}+(n-1)$ ②，①-②得，$2S_n+n^2-2S_{n-1}-(n-1)^2=2na_n+n-2(n-1)a_{n-1}-(n-1)$，即 $2a_n+2n-1=2na_n-2(n-1)a_{n-1}+1$，即 $2(n-1)a_n-2(n-1)a_{n-1}=2(n-1)$，所以 $a_n-a_{n-1}=1$, $n\ge 2$ 且 $n\in\mathbb{N}^*$，所以 $\{a_n\}$ 是以1为公差的等差数列.
(2) 由(1)可得 $a_4=a_1+3$, $a_7=a_1+6$, $a_9=a_1+8$，又 $a_4, a_7, a_9$ 成等比数列，所以 $a_7^2=a_4\cdot a_9$，即 $(a_1+6)^2=(a_1+3)(a_1+8)$，解得 $a_1=-12$，所以 $a_n=n-13$，所以 $S_n=-12n+\frac{n(n-1)}{2}=\frac{1}{2}n^2-\frac{25}{2}n=\frac{1}{2}\left(n-\frac{25}{2}\right)^2-\frac{625}{8}$，所以当 $n=12$ 或 $n=13$ 时，$(S_n)_{\min}=-78$.
\end{solution}
\end{question}

\begin{question}[points=12]
在四棱锥 $P-ABCD$ 中，$PD\perp$ 底面 $ABCD$, $CD\parallel AB$, $AD=DC=CB=1$, $AB=2$, $DP=\sqrt{3}$.

(1) 证明：$BD\perp PA$;

(2) 求 $PD$ 与平面 $PAB$ 所成的角的正弦值.

\begin{center}
  \includegraphics[width=.38\linewidth]{question-18-1.jpg}
\end{center}
\begin{solution}
\textbf{答案：}(1) 证明见解析；(2) $\frac{\sqrt{5}}{5}$

\textbf{分析：}(1) 作 $DE\perp AB$ 于E，$CF\perp AB$ 于F，利用勾股定理证明 $AD\perp BD$，根据线面垂直的性质可得 $PD\perp BD$，从而可得 $BD\perp$ 平面 $PAD$，再根据线面垂直的性质即可得证. (2) 以点D为原点建立空间直角坐标系，利用向量法即可得出答案.

\textbf{详解：}(1) 在四边形 $ABCD$ 中，作 $DE\perp AB$ 于E，$CF\perp AB$ 于F，因为 $CD\parallel AB$, $AD=CD=CB=1$, $AB=2$，所以四边形 $ABCD$ 为等腰梯形，所以 $AE=BF=\frac{1}{2}$，故 $DE=\frac{\sqrt{3}}{2}$, $BD=\sqrt{DE^2+BE^2}=\sqrt{3}$，所以 $AD^2+BD^2=AB^2$，所以 $AD\perp BD$. 因为 $PD\perp$ 平面 $ABCD$, $BD\subset$ 平面 $ABCD$，所以 $PD\perp BD$，又 $PD\cap AD=D$，所以 $BD\perp$ 平面 $PAD$，又因为 $PA\subset$ 平面 $PAD$，所以 $BD\perp PA$.
(2) 以点D为原点建立空间直角坐标系，$BD=\sqrt{3}$，则 $A(1,0,0)$, $B(0,\sqrt{3},0)$, $P(0,0,\sqrt{3})$，则 $\overrightarrow{AP}=(-1,0,\sqrt{3})$, $\overrightarrow{BP}=(0,-\sqrt{3},\sqrt{3})$, $\overrightarrow{DP}=(0,0,\sqrt{3})$，设平面 $PAB$ 的法向量 $\vec{n}=(x,y,z)$，则有 $\begin{cases} \vec{n}\cdot\overrightarrow{AP}=-x+\sqrt{3}z=0 \\ \vec{n}\cdot\overrightarrow{BP}=-\sqrt{3}y+\sqrt{3}z=0 \end{cases}$，可取 $\vec{n}=(\sqrt{3},1,1)$，则 $\cos\langle\vec{n},\overrightarrow{DP}\rangle=\frac{\vec{n}\cdot\overrightarrow{DP}}{|\vec{n}||\overrightarrow{DP}|}=\frac{\sqrt{3}}{\sqrt{5}\times\sqrt{3}}=\frac{\sqrt{5}}{5}$，所以 $PD$ 与平面 $PAB$ 所成角的正弦值为 $\frac{\sqrt{5}}{5}$.
\end{solution}
\end{question}

\begin{question}[points=12]
甲、乙两个学校进行体育比赛，比赛共设三个项目，每个项目胜方得10分，负方得0分，没有平局.三个项目比赛结束后，总得分高的学校获得冠军.已知甲学校在三个项目中获胜的概率分别为0.5, 0.4, 0.8，各项目的比赛结果相互独立.

(1) 求甲学校获得冠军的概率;

(2) 用X表示乙学校的总得分，求X的分布列与期望.
\begin{solution}
\textbf{答案：}(1) $0.6$；(2) 分布列见解析，$E(X)=13$

\textbf{分析：}(1) 设甲在三个项目中获胜的事件依次记为A,B,C，再根据甲获得冠军则至少获胜两个项目，利用互斥事件的概率加法公式以及相互独立事件的乘法公式即可求出. (2) 依题可知，X的可能取值为0,10,20,30，再分别计算出对应的概率，列出分布列，即可求出期望.

\textbf{详解：}(1) 设甲在三个项目中获胜的事件依次记为A,B,C，所以甲学校获得冠军的概率为 $P=P(ABC)+P(A\bar{B}C)+P(\bar{A}BC)+P(AB\bar{C})=0.5\times0.4\times0.8+0.5\times0.4\times0.8+0.5\times0.6\times0.8+0.5\times0.4\times0.2=0.16+0.16+0.24+0.04=0.6$.
(2) 依题可知，X的可能取值为0,10,20,30，所以 $P(X=0)=0.5\times0.4\times0.8=0.16$, $P(X=10)=0.5\times0.4\times0.8+0.5\times0.6\times0.8+0.5\times0.4\times0.2=0.44$, $P(X=20)=0.5\times0.6\times0.8+0.5\times0.4\times0.2+0.5\times0.6\times0.2=0.34$, $P(X=30)=0.5\times0.6\times0.2=0.06$. 即X的分布列为

| X | 0 | 10 | 20 | 30 |
|---|---|----|----|----|
| P | 0.16 | 0.44 | 0.34 | 0.06 |

期望 $E(X)=0\times0.16+10\times0.44+20\times0.34+30\times0.06=13$.
\end{solution}
\end{question}

\begin{question}[points=12]
设抛物线 $C:y^2=2px\ (p>0)$ 的焦点为F，点 $D(p,0)$，过F的直线交C于M，N两点.当直线MD垂直于x轴时，$|MF|=3$.

(1) 求C的方程;

(2) 设直线MD, ND与C的另一个交点分别为A，B，记直线MN, AB的倾斜角分别为 $\alpha$, $\beta$.当 $\alpha-\beta$ 取得最大值时，求直线AB的方程.
\begin{solution}
\textbf{答案：}(1) $y^2=4x$；(2) $AB: x=\sqrt{2}y+4$

\textbf{分析：}(1) 由抛物线的定义可得 $MF=p+\frac{p}{2}=3$，即可得解. (2) 设点的坐标及直线 $MN: x=my+1$，由韦达定理及斜率公式可得 $k_{MN}=2k_{AB}$，再由差角的正切公式及基本不等式可得 $k_{AB}=\frac{\sqrt{2}}{2}$，设直线 $AB: x=\sqrt{2}y+n$，结合韦达定理可解.

\textbf{详解：}(1) 抛物线的准线为 $x=-\frac{p}{2}$，当MD与x轴垂直时，点M的横坐标为p，此时 $|MF|=p+\frac{p}{2}=3$，所以 $p=2$，所以抛物线C的方程为 $y^2=4x$.
(2) 设 $M(\frac{y_1^2}{4},y_1)$, $N(\frac{y_2^2}{4},y_2)$, $A(\frac{y_3^2}{4},y_3)$, $B(\frac{y_4^2}{4},y_4)$，直线 $MN:x=my+1$，由 $\begin{cases} x=my+1 \\ y^2=4x \end{cases}$ 得 $y^2-4my-4=0$, $\Delta>0$, $y_1y_2=-4$. 由斜率公式可得 $k_{MN}=\frac{y_1-y_2}{\frac{y_1^2}{4}-\frac{y_2^2}{4}}=\frac{4}{y_1+y_2}$, $k_{AB}=\frac{y_3-y_4}{\frac{y_3^2}{4}-\frac{y_4^2}{4}}=\frac{4}{y_3+y_4}$. 直线 $MD:x=\frac{x_1-2}{y_1}\cdot y+2$，代入抛物线方程可得 $y^2-\frac{4(x_1-2)}{y_1}y-8=0$, $\Delta>0$, $y_1y_3=-8$，所以 $y_3=2y_2$，同理可得 $y_4=2y_1$，所以 $k_{AB}=\frac{4}{y_3+y_4}=\frac{4}{2(y_1+y_2)}=\frac{k_{MN}}{2}$. 又因为直线MN、AB的倾斜角分别为 $\alpha,\beta$，所以 $k_{AB}=\tan\beta=\frac{k_{MN}}{2}=\frac{\tan\alpha}{2}$，若要使 $\alpha-\beta$ 最大，则 $\beta\in(0,\frac{\pi}{2})$，设 $k_{MN}=2k_{AB}=2k>0$，则 $\tan(\alpha-\beta)=\frac{\tan\alpha-\tan\beta}{1+\tan\alpha\tan\beta}=\frac{k}{1+2k^2}=\frac{1}{\frac{1}{k}+2k}\leq\frac{1}{2\sqrt{\frac{1}{k}\cdot2k}}=\frac{\sqrt{2}}{4}$，当且仅当 $\frac{1}{k}=2k$ 即 $k=\frac{\sqrt{2}}{2}$ 时等号成立，所以当 $\alpha-\beta$ 最大时，$k_{AB}=\frac{\sqrt{2}}{2}$，设直线 $AB: x=\sqrt{2}y+n$，代入抛物线方程可得 $y^2-4\sqrt{2}y-4n=0$, $\Delta>0$, $y_3y_4=-4n=4y_1y_2=-16$，所以 $n=4$，所以直线 $AB: x=\sqrt{2}y+4$.
\end{solution}
\end{question}

\begin{question}[points=12]
已知函数 $f(x)=\frac{e^x}{x}-\ln x+x-a$.

(1) 若 $f(x)\geq 0$，求a的取值范围;

(2) 证明：若 $f(x)$ 有两个零点 $x_1, x_2$，则 $x_1x_2<1$.
\begin{solution}
\textbf{答案：}(1) $(-\infty, e+1]$；(2) 证明见解析

\textbf{分析：}(1) 由导数确定函数单调性及最值，即可得解. (2) 利用分析法，转化要证明条件为 $\frac{e^x}{x}-xe^{\frac{1}{x}}-2\left[\ln x-\frac{1}{2}\left(x-\frac{1}{x}\right)\right]>0$，再利用导数即可得证.

\textbf{详解：}(1) $f(x)$ 的定义域为 $(0,+\infty)$，则 $f'(x)=\left(\frac{1}{x}-\frac{1}{x^2}\right)e^x-\frac{1}{x}+1=\frac{x-1}{x}\left(\frac{e^x}{x}+1\right)$，令 $f'(x)=0$ 得 $x=1$，当 $x\in(0,1)$ 时 $f'(x)<0$, $f(x)$ 单调递减，当 $x\in(1,+\infty)$ 时 $f'(x)>0$, $f(x)$ 单调递增，$f(x)\geq f(1)=e+1-a$，若 $f(x)\geq0$，则 $e+1-a\geq0$，即 $a\leq e+1$，所以a的取值范围为 $(-\infty, e+1]$.
(2) 由题知，$f(x)$ 一个零点小于1，一个零点大于1，不妨设 $x_1<1<x_2$，要证 $x_1x_2<1$，即证 $x_1<\frac{1}{x_2}$，因为 $x_1,\frac{1}{x_2}\in(0,1)$，即证 $f(x_1)>f\left(\frac{1}{x_2}\right)$，又因为 $f(x_1)=f(x_2)$，故只需证 $f(x_2)>f\left(\frac{1}{x_2}\right)$，即证 $\frac{e^{x_2}}{x_2}-\ln x_2+x_2-\left(\frac{e^{\frac{1}{x_2}}}{\frac{1}{x_2}}-\ln\frac{1}{x_2}+\frac{1}{x_2}\right)>0$，即证 $\frac{e^{x_2}}{x_2}-x_2 e^{\frac{1}{x_2}}-2\left[\ln x_2-\frac{1}{2}\left(x_2-\frac{1}{x_2}\right)\right]>0$, $x_2>1$. 下面证明 $x>1$ 时，$\frac{e^x}{x}-xe^{\frac{1}{x}}>0$ 且 $\ln x-\frac{1}{2}\left(x-\frac{1}{x}\right)<0$. 设 $g(x)=\frac{e^x}{x}-xe^{\frac{1}{x}}$, $x>1$，则 $g'(x)=\left(\frac{1}{x}-\frac{1}{x^2}\right)e^x-\left(e^{\frac{1}{x}}+xe^{\frac{1}{x}}\cdot(-\frac{1}{x^2})\right)=\frac{x-1}{x}\left(\frac{e^x}{x}-e^{\frac{1}{x}}\right)$，设 $\varphi(x)=\frac{e^x}{x}\ (x>1)$，$\varphi'(x)=\frac{x-1}{x^2}e^x>0$，所以 $\varphi(x)>\varphi(1)=e$，而 $e^{\frac{1}{x}}<e$，所以 $\frac{e^x}{x}-e^{\frac{1}{x}}>0$，故 $g'(x)>0$，$g(x)$ 在 $(1,+\infty)$ 单调递增，$g(x)>g(1)=0$，所以 $\frac{e^x}{x}-xe^{\frac{1}{x}}>0$. 令 $h(x)=\ln x-\frac{1}{2}\left(x-\frac{1}{x}\right)$, $x>1$，则 $h'(x)=\frac{1}{x}-\frac{1}{2}\left(1+\frac{1}{x^2}\right)=\frac{2x-x^2-1}{2x^2}=\frac{-(x-1)^2}{2x^2}<0$，所以 $h(x)$ 在 $(1,+\infty)$ 单调递减，$h(x)<h(1)=0$，所以 $\ln x-\frac{1}{2}\left(x-\frac{1}{x}\right)<0$. 综上，$\frac{e^x}{x}-xe^{\frac{1}{x}}-2\left[\ln x-\frac{1}{2}\left(x-\frac{1}{x}\right)\right]>0$，所以 $x_1x_2<1$.
\end{solution}
\end{question}

\begin{question}[points=10]
在直角坐标系 $xOy$ 中，曲线 $C_1$ 的参数方程为 $\begin{cases} x=\frac{2+t}{6} \\ y=\sqrt{t} \end{cases}$ (t为参数)，曲线 $C_2$ 的参数方程为 $\begin{cases} x=-\frac{2+s}{6} \\ y=-\sqrt{s} \end{cases}$ (s为参数).

(1) 写出 $C_1$ 的普通方程;

(2) 以坐标原点为极点，x轴正半轴为极轴建立极坐标系，曲线 $C_3$ 的极坐标方程为 $2\cos\theta-\sin\theta=0$，求 $C_3$ 与 $C_1$ 交点的直角坐标，及 $C_3$ 与 $C_2$ 交点的直角坐标.
\begin{solution}
\textbf{答案：}(1) $y^2=6x-2\ (y\geq0)$；(2) $C_3,C_1$ 的交点坐标为 $(\frac{1}{2},1)$ 和 $(1,2)$；$C_3,C_2$ 的交点坐标为 $(-\frac{1}{2},-1)$ 和 $(-1,-2)$

\textbf{分析：}(1) 消去t，即可得到 $C_1$ 的普通方程. (2) 将曲线 $C_2, C_3$ 的方程化成普通方程，联立求解即可.

\textbf{详解：}(1) 因为 $x=\frac{2+t}{6}$, $y=\sqrt{t}$，所以 $x=\frac{2+y^2}{6}$，即 $C_1$ 的普通方程为 $y^2=6x-2\ (y\geq0)$.
(2) 因为 $x=-\frac{2+s}{6}$, $y=-\sqrt{s}$，所以 $6x=-2-y^2$，即 $C_2$ 的普通方程为 $y^2=-6x-2\ (y\leq0)$. 由 $2\cos\theta-\sin\theta=0$ 得 $2\rho\cos\theta-\rho\sin\theta=0$，即 $C_3$ 的普通方程为 $2x-y=0$. 联立 $\begin{cases} y^2=6x-2\ (y\geq0) \\ 2x-y=0 \end{cases}$，解得 $\begin{cases} x=\frac{1}{2} \\ y=1 \end{cases}$ 或 $\begin{cases} x=1 \\ y=2 \end{cases}$，即交点坐标为 $(\frac{1}{2},1)$ 和 $(1,2)$；联立 $\begin{cases} y^2=-6x-2\ (y\leq0) \\ 2x-y=0 \end{cases}$，解得 $\begin{cases} x=-\frac{1}{2} \\ y=-1 \end{cases}$ 或 $\begin{cases} x=-1 \\ y=-2 \end{cases}$，即交点坐标为 $(-\frac{1}{2},-1)$ 和 $(-1,-2)$.
\end{solution}
\end{question}

\begin{question}[points=10]
已知a, b, c均为正数，且 $a^2+b^2+4c^2=3$，证明:

(1) $a+b+2c\leq 3$;

(2) 若 $b=2c$，则 $\frac{1}{a}+\frac{1}{c}\geq 3$.
\begin{solution}
\textbf{答案：}(1) 证明见解析；(2) 证明见解析

\textbf{分析：}(1) 方法一：根据 $a^2+b^2+4c^2=a^2+b^2+(2c)^2$，利用柯西不等式即可得证. (2) 由(1)结合已知可得 $0<a+4c\leq 3$，即可得到 $\frac{1}{a+4c}\geq \frac{1}{3}$，再根据权方和不等式即可得证.

\textbf{详解：}(1) 由柯西不等式有 $[a^2+b^2+(2c)^2](1^2+1^2+1^2)\geq (a+b+2c)^2$，所以 $a+b+2c\leq 3$，当且仅当 $a=b=2c=1$ 时取等号，所以 $a+b+2c\leq 3$.
(2) 因为 $b=2c$, $a>0$, $b>0$, $c>0$，由(1)得 $a+b+2c=a+4c\leq 3$，即 $0<a+4c\leq 3$，所以 $\frac{1}{a+4c}\geq \frac{1}{3}$. 由权方和不等式知 $\frac{1}{a}+\frac{1}{c}=\frac{1}{a}+\frac{2}{2c}\geq \frac{(1+\sqrt{2})^2}{a+2c}=\frac{(1+\sqrt{2})^2}{a+\frac{b}{2}}$? 注意 $b=2c$，则 $c=\frac{b}{2}$，且 $b=2c$，所以 $c=\frac{b}{2}$，而 $b$ 是变量，这里需改用另一形式：$\frac{1}{a}+\frac{1}{c}=\frac{1^2}{a}+\frac{1^2}{c}\geq \frac{(1+1)^2}{a+c}=\frac{4}{a+c}$，但 $a+c$ 与 $a+4c$ 不同，需重新推导：由 $b=2c$ 及条件 $a^2+b^2+4c^2=3$ 得 $a^2+4c^2+4c^2=3$，即 $a^2+8c^2=3$，且 $a+4c\leq 3$，但目标为 $\frac{1}{a}+\frac{1}{c}\geq 3$，可考虑用柯西不等式：$(\frac{1}{a}+\frac{1}{c})(a+4c)=(\frac{1}{a}+\frac{1}{c})(a+4c)\geq (1+2)^2=9$，所以 $\frac{1}{a}+\frac{1}{c}\geq \frac{9}{a+4c}\geq \frac{9}{3}=3$，当且仅当 $\frac{1}{a}:\frac{1}{c}=a:4c$ 即 $a=2c$ 且 $b=2c$ 时取等号，此时 $a=2c$, $a^2+8c^2=4c^2+8c^2=12c^2=3$，得 $c=\frac{1}{2}$, $a=1$，等号成立. 故 $\frac{1}{a}+\frac{1}{c}\geq 3$.
\end{solution}
\end{question}

